\documentclass[12pt]{article}
\usepackage{amsfonts,amsthm,amsmath,amssymb,mathtools}
\usepackage{mathrsfs}
\usepackage{enumitem}
\usepackage{tikz-cd}
\usepackage{geometry}
\usepackage{graphicx}
\IfFileExists{mlmodern.sty}{
    \usepackage{mlmodern}
}{}

\geometry{letterpaper, portrait, margin=1in}

\setcounter{section}{0}

\renewcommand{\thesection}{\S\arabic{section}}
\renewcommand{\thesubsection}{\arabic{section}}
\DeclareMathOperator{\im}{im}
\DeclareMathOperator{\id}{id}
\DeclareMathOperator{\ad}{ad}
\newcommand{\gl}{\mathfrak{gl}}
\renewcommand{\sl}{\mathfrak{sl}}
\renewcommand{\th}{^\mathrm{th}}

\newtheorem{thm}{Theorem}
\newenvironment{theorem}[1]{
	\renewcommand{\thethm}{#1}
	\begin{thm}
		}{\end{thm}}

\newtheorem{lma}{Lemma}
\newenvironment{lemma}[1]{
	\renewcommand{\thelma}{#1}
	\begin{lma}
		}{\end{lma}}

\theoremstyle{definition}
\newtheorem{ex}{}
\newenvironment{exercise}[1]{
	\renewcommand{\theex}{#1}
	\begin{ex}
		}{\end{ex}}

\title{Lie Theory: Chapter 2 Homework}
\author{Connor Mooneyhan}
\date{}

\begin{document}

\maketitle

\begin{ex}
	Let $I$ and $J$ be ideals of a Lie algebra $L$.
	Prove that $I \cap J$ is a subspace of $L$.
\end{ex}
\begin{proof}
	Given $a, b \in F$ and $x, y \in I \cap J$, we want to show that $ax + by \in I\cap J$.
	Indeed, $ax+by \in I$ since $I$ is a subspace by virtue of being an ideal, and $ax + by \in J$ for the same reason, so $ax + by \in I \cap J$.
\end{proof}

\begin{ex}
	Prove the Isomorphism Theorem (part a).
\end{ex}
\begin{proof}
	Given Lie algebras $L_1$ and $L_2$ together with a homomorphism $f : L_1 \to L_2$, define $g : L_1/\ker f \to L_2$ by $g(x + \ker f) = f(x)$.
	We first show that this is well-defined.

	If $x + \ker f = y + \ker f$, we have $y - x \in \ker f$, so \begin{align*}
		g(x + \ker f) & = f(x)            \\
		              & = f(x) + 0        \\
		              & = f(x) + f(y - x) \\
		              & = f(x + y - x)    \\
		              & = f(y)            \\
		              & = g(y + \ker f),
	\end{align*}
	so $g$ is well-defined.

	We now show that it is a linear map by taking, for any $x, y \in L_1$, \begin{align*}
		g(a(x + \ker f) + b(y + \ker f)) & = g((ax + \ker f) + (by + \ker f)) \\
		                                 & = g((ax + by) + \ker f)            \\
		                                 & = f(ax + by)                       \\
		                                 & = af(x) + bf(y)                    \\
		                                 & = ag(x + \ker f) + bg(y + \ker f).
	\end{align*}
	It is furthermore a Lie algebra homomorphism because \begin{align*}
		g([x + \ker f, y + \ker f]) & = g([x, y] + \ker f)              \\
		                            & = f([x, y])                       \\
		                            & = [f(x), f(y)]                    \\
		                            & = [g(x + \ker f), g(y + \ker f)].
	\end{align*}

	Finally, we show bijectivity onto its image.
	It is surjective onto $\im g = \im f$ by definition.
	For injectivity, let $x, y \in L_1$ such that $g(x + \ker f) = g(y + \ker f)$.
	Then \begin{align*}
		0 & = g(x + \ker f) - g(y + \ker  f) \\
		  & = f(x) - f(y)                    \\
		  & = f(x - y),
	\end{align*}
	so $x - y \in \ker f$, from which it follows that $x + \ker f = y + \ker f$.
	Thus $f$ is a bijective Lie algebra homomorphism $L_1\ker f \to \im f$, i.e. an isomorphism.
\end{proof}

\begin{ex}
	Prove the statements on page 14 (just above example 2.4): \begin{enumerate}[label=(\alph*)]
		\item Prove that $J$ (as defined on page 14) is an ideal of $L$.
		\item Prove that $J$ contains $I$.
	\end{enumerate}
\end{ex}
\begin{proof}
	Let $K$ be an ideal of $L/I$ and define $J \coloneq \lbrace z \in L : z + I \in K \rbrace$.
	\begin{enumerate}[label=(\alph*)]
		\item Given $a, b \in F$ and $x, y \in J$, \begin{align*}
			      (ax + by) + I & = (ax + I) + (by + I)  \\
			                    & = a(x + I) + b(y + I),
		      \end{align*}
		      so since $x + I$ and $y + I$ are in $K$, this is a linear combination of elements of $K$, which is itself in $K$ as desired.
		      Thus $J$ is a linear subspace.

		      Now if $x \in J$ and $y \in L$, we have \begin{align*}
			      [x, y] + I & = [x + I, y + I],
		      \end{align*}
		      so since $x + I \in K$ and $K$ is an ideal, we get $[x, y] + I \in K$ and thus $[x, y] \in J$.
		\item Given $x \in I$, $x + I = I$, so since $I$ is the identity in $L/I$ and every ideal contains the identity, $x + I \in K$ and thus $x \in J$.
	\end{enumerate}
\end{proof}

\begin{ex}
	Suppose $L_1$ and $L_2$ are Lie algebras.
	Let $L \coloneq \lbrace (x_1, x_2) : x_i \in L_i \rbrace$ be the direct sum of their underlying vector spaces. \begin{enumerate}[label=(\alph*)]
		\item Show that if we define \begin{equation*}
			      [(x_1, x_2), (y_1, y_2)] \coloneq ([x_1, y_1], [x_2, y_2])
		      \end{equation*}
		      then $L$ becomes a Lie algebra, the \textit{direct sum} of $L_1$ and $L_2$.
		      We denote this by $L = L_1 \oplus L_2$.
		\item Show that if $L = L_1 \oplus L_2$ then $Z(L) = Z(L_1) \oplus Z(L_2)$ and $L' = L'_1 \oplus L'_2$.
	\end{enumerate}
\end{ex}
\begin{proof}
	\begin{enumerate}[label=(\alph*)]
		\item Since we already know that $L$ is a vector space, we will show that the bracket as defined is a valid Lie bracket.
		      Given $(x_1, x_2) \in L$, we have \begin{align*}
			      [(x_1, x_2), (x_1, x_2)] & = ([x_1, x_1], [x_2, x_2]) \\
			                               & = (0, 0)
		      \end{align*}
		      as desired.
		      For $(y_1, y_2)$ and $(z_1, z_2)$ also in $L$, we have \begin{align*}
			          & [(x_1, x_2), [(y_1, y_2), (z_1, z_2)]] + [(y_1, y_2), [(z_1, z_2), (x_1, x_2)]] + [(z_1, z_2), [(x_1, x_2), (y_1, y_2)]] \\
			      =\  & [(x_1, x_2), ([y_1, z_1], [y_2, z_2])] + [(y_1, y_2), ([z_1, x_1], [z_2, x_2])] + [(z_1, z_2), ([x_1, y_1], [x_2, y_2])] \\
			      =\  & ([x_1, [y_1, z_1]], [x_2, [y_2, z_2]]) + ([y_1, [z_1, x_1]], [y_2, [z_2, x_2]]) + ([z_1, [x_1, y_1]], [z_2, [x_2, y_2]]) \\
			      =\  & ([x_1, [y_1, z_1]] + [y_1, [z_1, x_1]] + [z_1, [x_1, y_1]], [x_2, [y_2, z_2]] + [y_2, [z_2, x_2]] + [z_2, [x_2, y_2]])   \\
			      =\  & (0, 0)
		      \end{align*}
		      as desired.
		\item If $(x_1, x_2) \in Z(L)$, then given any $(y_1, y_2) \in L$, \begin{equation*}
			      (0, 0) = [(x_1, x_2), (y_1, y_2)] = ([x_1, y_1], [x_2, y_2]),
		      \end{equation*}
		      so $[x_1, y_1] = [x_2, y_2] = 0$, and since $y_1$ and $y_2$ were arbitrary, this means $(x_1, x_2) \in Z(L_1) \oplus Z(L_2)$.
		      Also, if $(x_1, x_2) \in Z(L_1) \oplus Z(L_2)$, then for any $(y_1, y_2) \in L$, \begin{equation*}
			      [(x_1, x_2), (y_1, y_2)] = ([x_1, y_1], [x_2, y_2]) = (0, 0),
		      \end{equation*}
		      so $(x_1, x_2) \in Z(L)$.
		      Thus by mutual inclusion, $Z(L) = Z(L_1) \oplus Z(L_2)$.

		      Now if $x \in L'$, then $x = \sum_{i < j} a_{ij}[x^i, x^j]$ for some $\lbrace x^1, x^2, \dots, x^n \rbrace \subset L$ where the superscripts are to be interpreted just as indexing notation, not exponents, and where each $a_{ij} \in F$.
		      We have, \begin{align*}
			      x & = \sum_{i < j} a_{ij}[x^i, x^j]                                                                         \\
			        & = \sum_{i < j} a_{ij}[(x^i_1, x^i_2), (x^j_1, x^j_2)]                                                   \\
			        & = \sum_{i < j} a_{ij}([x^i_1, x^j_1], [x^i_2, x^j_2])                                                   \\
			        & = \sum_{i < j} (a_{ij}[x^i_1, x^j_1], a_{ij}[x^i_2, x^j_2])                                             \\
			        & = \left(\sum_{i < j}a_{ij}[x^i_1, x^j_1],\sum_{i < j} a_{ij}[x^i_2, x^j_2]\right) \in L_1' \oplus L_2'.
		      \end{align*}
		      Running that argument backwards, starting with an element of $L_1' \oplus L_2'$, shows the reverse containment, so we have $L' = L_1' \oplus L_2'$.
	\end{enumerate}
\end{proof}

\begin{ex}
	For each pair of the following Lie algebras over $\mathbb{R}$, decide whether or not they are isomorphic: \begin{enumerate}[label=(\alph*)]
		\item the Lie algebra $\mathbb{R}^3_\wedge$ with Lie bracket given by the vector product;
		\item the upper triangular $2 \times 2$ matrices over $\mathbb{R}$;
		\item the strictly upper triangular $3 \times 3$ matrices over $\mathbb{R}$;
		\item $L = \lbrace x \in \gl(3, \mathbb{R}) : x^t = -x \rbrace$.
	\end{enumerate}
\end{ex}
\begin{proof}
	We note a few things up front.
	All four candidates have dimension 3.
	We now look at the dimension of the derived algebra for each. \begin{equation*}
		\begin{array}{r|l}
			L                        & \mathrm{dim}(L')                                                                                                                      \\\hline
			\mathbb{R}^3_\wedge      & \text{3; every vector in $\mathbb{R}^3$ can be expressed as the}                                                                      \\
			                         & \text{product of two others.}                                                                                                          \\
			b(2, \mathbb{R})         & \text{1; the bracket of two upper triangular matrices is \textit{strictly} upper triangular,}                                         \\
			                         & \text{as we have seen, leaving us with nonzero entries in the top-right entry only.}                                                   \\
			n(3, \mathbb{R})         & \text{1; similar to the above case, where the bracket leaves us with nonzero entries}                                                 \\
			                         & \text    {in the top-right entry only.}                                                                                                \\
			\lbrace x^t = -x \rbrace & \text{3; the matrices here look like $\begin{pmatrix} 0 & a & b \\ -a & 0 & c \\ -b & -c & 0\end{pmatrix}$, so there are three basis} \\
        & \text{elements, each taking $a = 1$, $b = 1$, and $c = 1$ respectively, with zeroes}                                                   \\
                               & \text{everywhere else.}
		\end{array}
	\end{equation*}
  So the only candidates for isomorphism are (a) with (d) and (b) with (c).
  We also have that (c) is associative, whereas (b) is nonassociative, so (c) is not isomorphic to (b).

	\textcolor{red}{UNFINISHED}
\end{proof}

\begin{ex}
	Let $A$ and $B$ be $n \times n$ matrices with complex entries.
	\begin{enumerate}[label=(\alph*)]
		\item Is it true that $tr(AB) = tr(A)tr(B)$?
		      If true, prove it.
		      If false, provide a counterexample.
		\item Is it true that $tr(AB) = tr(BA)$?
		      If true, prove it.
		      If false, provide a counterexample.
		\item Use your answers to (a) and/or (b) to show that if $x \in \gl_S(n, \mathbb{C})$ for some $n \times n$ invertible matrix $S$, then $tr(x) = 0$.
	\end{enumerate}
\end{ex}
\begin{proof}\
	\begin{enumerate}[label=(\alph*)]
		\item No.
		      Consider \begin{equation*}
			      tr \left( \begin{pmatrix}1 & 0 \\ 0 & 1\end{pmatrix}\begin{pmatrix}1 & 0 \\ 0 & 1\end{pmatrix} \right) = tr \left( \begin{pmatrix}1 & 0 \\ 0 & 1\end{pmatrix} \right) = 2
		      \end{equation*}
		      and \begin{equation*}
			      tr \left( \begin{pmatrix}1 & 0 \\ 0 & 1\end{pmatrix}\right) tr\left(\begin{pmatrix}1 & 0 \\ 0 & 1\end{pmatrix} \right) = 2 \cdot 2 = 4.
		      \end{equation*}
		\item Yes. \begin{align*}
			      tr(AB) & = \sum_{i=1}^n (AB)_{ii}                 \\
			             & = \sum_{i=1}^n \sum_{j=1}^n A_{ij}B_{ji} \\
			             & = \sum_{i=1}^n \sum_{j=1}^n A_{ji}B_{ij} \\
			             & = \sum_{i=1}^n \sum_{j=1}^n B_{ij}A_{ji} \\
			             & = \sum_{i=1}^n (BA)_{ii}                 \\
			             & = tr(BA),
		      \end{align*}
		      where the third equality comes not from term-by-term equality as ordered, but from realizing that since we are summing over all $i$ and $j$ between 1 and $n$ inclusive, each term from the sum is represented once in both versions, so we can use commutativity to reorder the sum conveniently to fit our needs.
		\item If $x \in \gl_S(n, \mathbb{C})$, then $x^TS = -Sx$.
		      We will also use the fact that $tr(y) = tr(y^T)$ for any matrix $y$ since the diagonal doesn't change when you take the transpose.
		      We have \begin{align*}
			      tr(x) & = tr(S^{-1}Sx)      \\
			            & = -tr(S^{-1}(x^TS)) \\
			            & = -tr(x^TSS^{-1})   \\
			            & = -tr(x^T)          \\
			            & = -tr(x),
		      \end{align*}
		      so $tr(x) = 0$.
	\end{enumerate}
\end{proof}

\begin{ex}
	Let $I$ be an ideal of a Lie algebra $L$.
	Let $B$ be the centralizer of $I$ in $L$; that is \begin{equation*}
		B = C_L(I) = \left\lbrace x \in L : [x, a] = 0 \text{ for all } a \in I \right\rbrace.
	\end{equation*}
	Show that $B$ is an ideal of $L$.
\end{ex}
\begin{proof}
	First, we show closure under linear combinations.
	If $a, b \in F$, $x, y \in B$, and $z \in I$, \begin{align*}
		[ax + by, z] & = a[x, z] + b[y, z] \\
		             & = 0 + 0             \\
		             & = 0,
	\end{align*}
	so $ax + by \in B$.

	Now if $x \in B$ and $y \in L$, then given any $z \in I$ \begin{align*}
		[[x, y], z] & = -[z, [x, y]]               \\
		            & = [x, [y, z]] + [y, [z, x]].
	\end{align*}
	Note that $[y, z] \in I$ since $z \in I$, so $[x, [y, z]] = 0$, and $[z, x] = 0$ since $x \in B$ so $[y, [z, x]] = [y, 0] = 0$, so we indeed have $[[x,y],z] = 0$, meaning $[x,y] \in B$.
	Thus $B$ is an ideal.
\end{proof}

\begin{ex}
	Prove the Isomorphism Theorem (part b).
\end{ex}
\begin{proof}
	Define $f : I + J \to \frac{I}{I \cap J}$ by $i + j \mapsto i + I \cap J$.
	This is well defined because if $i_1 + j_1 = i_2 + j_2$ then $i_2 - i_1 = j_1 - j_2 \in I \cap J$, so \begin{align*}
		f(i_1 + j_1) & = i_1 + I \cap J               \\
		             & = i_1 + (i_2 - i_1 + I \cap J) \\
		             & = (i_1 + i_2 - i_1) + I \cap J \\
		             & = i_2 + I \cap J               \\
		             & = f(i_2 + j_2).
	\end{align*}
	It is linear because \begin{align*}
		f(a(i_1 + j_1) + b(i_2 + j_2)) & = f(ai_1 + bi_2 + aj_1 + bj_2)          \\
		                               & = ai_1 + bi_2 + I \cap J                \\
		                               & = (ai_1 + I \cap J) + (bi_2 + I \cap J) \\
		                               & = a(i_1 + I \cap J) + b(i_2 + I \cap J) \\
		                               & = af(i_1 + j_1) + bf(i_2 + j_2),
	\end{align*}
	and \begin{align*}
		f([i_1 + j_1, i_2 + j_2]) & = f([i_1, i_2] + [i_1, j_2] + [j_1, i_2] + [j_1, j_2]) \\
		                          & = [i_1, i_2] + I \cap J                                \\
		                          & = [i_1 + I \cap J, i_2 + I \cap J]                     \\
		                          & = [f(i_1 + j_1), f(i_2 + j_2)],
	\end{align*}
	where the second equality results from the fact that the final three terms of the argument of the RHS of the first equality are in $J$ by virtue of $J$ being an ideal.
	Thus $f$ is a Lie algebra homomorphism.

	We will now show that $\ker f = J$.
	Indeed, if $i \in I$ is not in $J$, then $i$ is not in $I \cap J$, so $f(i + j) \neq I \cap J$, i.e. it is nonzero, for any choice of $j \in J$.
	If $i \in I$ \textit{is} in $J$, then $i \in I \cap J$, so $f(i + j) = i + I \cap J = I \cap J$, which is zero in $\frac{I}{I \cap J}$.
	Clearly if $i + j$ is of the form $0 + j$, then $f(0 + j) = 0 + I \cap J = I \cap J$.
	Thus $\ker f = J$.

	We also see immediately that $f$ is surjective, since given $i + I \cap J \in \frac{I}{I \cap J}$, we have $f(i) = i + I \cap J$.
	Thus the first isomorphism theorem gives us that \begin{equation*}
		\frac{I + J}{J} = \frac{I + J}{\ker f} \cong \im f = \frac{I}{I \cap J}.
	\end{equation*}
\end{proof}

\begin{ex}
	Prove the Isomorphism Theorem (part c).
\end{ex}
\begin{proof}
	Let $I$ and $J$ be ideals of a Lie algebra $L$ such that $I \subset J$.

	We first show that $J / I$ is an ideal of $L / I$.
	Indeed, given $j + I \in J / I$ and $\ell + I \in L / I$, we have \begin{align*}
		[j + I, \ell + I] & = [j, \ell] + I \in J / I
	\end{align*}
	since $[j, \ell] \in J$ because $J$ is an ideal.

	Now we define $f : L/I \to L/J$ by $\ell + I \mapsto \ell + J$, and we want to show that this is surjective homomorphism and that its kernel is $J/I$.
	First we show that $f$ is well-defined: if $\ell_1 + I = \ell_2 + I$ then $\ell_2 - \ell_1 \in I \subset J$, so \begin{align*}
		f(\ell_1 + I) & = \ell_1 + J                     \\
		              & = \ell_1 + (\ell_2 - \ell_1 + J) \\
		              & = (\ell_1 + \ell_2 - \ell_1) + J \\
		              & = \ell_2 + J                     \\
		              & = f(\ell_2 + I).
	\end{align*}
	Now we show that it is linear: \begin{align*}
		f(a(\ell_1 + I) + b(\ell_2 + I)) & = f((a\ell_1 + b\ell_2) + I)       \\
		                                 & = (a\ell_1 + b\ell_2) + J          \\
		                                 & = (a\ell_1 + J) + (b\ell_2 + J)    \\
		                                 & = a(\ell_1 + J) + b(\ell_2 + J)    \\
		                                 & = af(\ell_1 + I) + bf(\ell_2 + I).
	\end{align*}
	Then we show that it preserves brackets: \begin{align*}
		f([\ell_1 + I, \ell_2 + I]) & = f([\ell_1, \ell_2] + I)         \\
		                            & = [\ell_1, \ell_2] + J            \\
		                            & = [\ell_1 + J, \ell_2 + J]        \\
		                            & = [f(\ell_1 + I), f(\ell_2 + I)].
	\end{align*}
	So $f$ is a Lie algebra homomorphism.

	Now, given any $\ell + J \in L/J$, $f(\ell + I) = \ell + J$, so we have surjectivity.
	Finally, if $f(\ell + I) = \ell + J = J$, then we must have $\ell \in J$, so $\ker f \subset J/I$.
	On the other hand, if $j \in J$, then $f(j + I) = j + J = J$, so we have $\ker f = J/I$ as desired.

	Thus the first isomorphism theorem gives us that \begin{equation*}
		\frac{L/I}{J/I} = \frac{L/I}{\ker f} \cong \im f = \frac{L}{J}.
	\end{equation*}
\end{proof}

\begin{ex}
	Let $F$ be a field.
	Show that the derived algebra of $\gl(n, F)$ is $\sl(n, F)$.
\end{ex}
\begin{proof}
	Let $A, B \in \gl(n, F)$.
	Since $tr(AB) = tr(BA)$ and trace is linear, \begin{align*}
		tr([A, B]) & = tr(AB - BA)     \\
		           & = tr(AB) - tr(BA) \\
		           & = tr(AB) - tr(AB) \\
		           & = 0.
	\end{align*}
	Thus $[A, B] \in \sl(n, F)$ for all $A, B \in \gl(n, F)$, so their span is contained in the subspace $\sl(n, F)$, i.e. $\gl(n, F)' \subset \sl(n, F)$.

	To show the other direction, we show that the standard basis elements of $\sl(n, F)$ are contained in $\gl(n, F)'$.
	We know that for $a,b,c,d$ between 1 and $n$, $e_{ab}e_{cd} = \delta_{bc}e_{ad}$ since all entries are zero except perhaps entry $ad$, which is nonzero if and only if $b = c$.
	Given this, we can express the basis elements of $\sl(n, F)$ as \begin{equation*}
		e_{ij} = e_{ii}e_{ij} = e_{ii}e_{ij} - e_{ij}e_{ii} = [e_{ii},e_{ij}] \in \gl(n, F)'
	\end{equation*} when $i \neq j$ and \begin{equation*}
		e_{ii} - e_{i+1,i+1} = e_{i, i + 1}e_{i+1, i} - e_{i+1,i}e_{i,i+1} = [e_{i,i+1}, e_{i+1,i}] \in \gl(n, F)'.
	\end{equation*}
	Thus $\sl(n, F) = \gl(n, F)'$ by mutual containment.
\end{proof}

\end{document}
